\documentclass{article}
\usepackage{xepersian}
\settextfont{XB Niloofar}

\begin{document}

\title{نمونه سند Mixed Document}
\author{Author Name نویسنده}
\date{\today}

\maketitle

\section{مقدمه Introduction}

این یک پاراگراف تستی است که شامل both Persian and English text می‌باشد. This makes it difficult to read در ادیتورهای مختلف. We want to separate این محتوا to make it more readable.

هنگامی که با AI می‌نویسیم، often we get mixed content که خواندن آن در ویرایشگرها مشکل است. This tool helps با جداسازی محتوای مخلوط.

\section{روش‌ها Methods}

این بخش شامل توضیحاتی درباره the methodology used در این تحقیق می‌باشد.

\subsection{نمونه‌ها Samples}

The samples were collected از منابع مختلف and processed using مدل‌های machine learning.

\section{نتایج Results}

نتایج نشان می‌دهد که the proposed method عملکرد بهتری نسبت به روش‌های قبلی دارد.

\begin{itemize}
    \item دقت accuracy بیش از 95 درصد
    \item زمان پردازش processing time کمتر از یک ثانیه
    \item مصرف حافظه memory usage بهینه شده
\end{itemize}

\section{نتیجه‌گیری Conclusion}

در این مقاله we presented یک روش جدید for processing mixed-language LaTeX documents که به way قابل توجهی readability را بهبود می‌بخشد.

این پاراگراف کاملاً به زبان فارسی نوشته شده است و نباید تغییر کند. همه کلمات در این پاراگراف فارسی هستند.

This paragraph is written entirely in English and should remain unchanged. All words in this paragraph are in English.

\end{document}
